\section{问题重述与分析}

\subsection{问题背景}

啤酒的香气与风味由醇、酯、酸、硫化物、萜烯等多类化合物共同决定。同一化合物在不同浓度区间可能呈现水果香、花香、麦芽香或异味，不同化合物之间还存在加成、掩盖与协同作用。仅凭单一化合物的浓度很难解释整体感官印象，因此需要把化学测量、感官评分与化合物之间的关联结构放在同一个框架中分析\cite{jiang2018}。本题给出 \textbf{250} 个啤酒样本的化学特征、\textbf{50} 个感官维度、化合物相关矩阵以及若干辅助属性表，要求依次完成感官评分预测、关键化合物识别和风格条件生成三项任务。数据驱动的建模方法为这类高维小样本问题提供了可行路径\cite{zhang2020}。

\subsection{问题一的重述与分析}

\textbf{重述：}问题一要求根据每个啤酒样本的化学特征，预测其在 \textbf{50} 个感官维度上的 \textbf{0--6} 级评分。训练数据包含 \textbf{231} 个化学特征和 \textbf{50} 个连续感官标签，标签需要按题面评分卡转成整数等级；评价时既关心连续预测的拟合程度，也关心分级结果是否保持原有顺序与分布形态。

\textbf{分析：}这一问的本质是“高维输入、小样本、多输出”的回归问题。难点有三：其一，\textbf{231} 个特征之间存在明显共线性，线性模型容易过拟合；其二，样本只有 \textbf{250} 个，而每个感官维度都要单独建模，折外预测的方差会被放大；其三，题面要求输出 \textbf{0--6} 的整数，而连续预测的分布通常比标签更集中，直接四舍五入会丢失两端等级。针对这些难点，本文用随机森林捕捉非线性与交互作用\cite{breiman2001}，用固定折索引的交叉验证控制评估偏差\cite{li2019}，再把连续预测送入预先冻结的 ALR 标定规则；对两端等级为空的现象，不通过拉伸预测去“补”，而是给出可检验的机理解释与可读口径辅助表。

\subsection{问题二的重述与分析}

\textbf{重述：}问题二要求利用化合物相关矩阵建立共表达网络，并结合化学特征与感官标签，为每个感官维度找出 \textbf{3} 个关键化合物；同时要说明结果的稳定性、化学合理性，并与全局重要度排序进行对照。

\textbf{分析：}这一问的核心不是“算一个相关系数”，而是“在相关结构上做可解释的排序”。直接按化合物与感官维度的相关系数排序，容易选中彼此高度共线的“同伙”，忽略与目标化合物协同起作用的邻居；只按网络中心性排序，又会脱离感官目标。因此本文采用“直接相关得分 + 邻居加权得分”的融合分数，把图结构作为平滑项而不是替代项。稳定性用有放回重采样衡量，化学合理性用化合物来源品类与感官维度的对应关系核对，全局对照只作为旁证，避免把恒真的集合包含关系误读为重要度一致。多重比较用 Benjamini--Hochberg 方法控制错误发现率\cite{benjamini1995}。

\subsection{问题三的重述与分析}

\textbf{重述：}问题三要求以啤酒风格为条件，学习一个能够生成化合物谱与对应感官画像的模型，并考察生成结果在不同风格之间的迁移能力。题面同时要求生成结果在化学上可解释、在感官上可核对。

\textbf{分析：}这一问与问题一、问题二的方向相反：前两问是“由化合物到感官”的判别与排序，这一问是“由风格到化合物谱”的条件生成。难点在于风格标签的样本量不均衡、化合物维度高且存在大量零值与强相关，生成器很容易只记住训练风格而失去迁移能力。本文拟采用“风格条件编码—化合物谱生成—感官回评—化学约束”的路线，把问题一的预测器作为回评器，把问题二的关键化合物作为化学约束的落点；本阶段只登记路线与验证设计，不报告任何生成结果。

\subsection{总体思路}

全文按“先统一口径、再分别建模、最后交叉检验”的顺序展开。第 \textbf{3} 章登记数据口径与预处理；第 \textbf{4} 章解决问题一，重点是随机森林回归、ALR 标定和刻度不可达命题；第 \textbf{5} 章解决问题二，重点是共表达图、融合排序、稳定性与化学合理性；第 \textbf{6} 章给出问题三的题面要求与拟采用路线；第 \textbf{7} 章评价模型并给出结论与建议。三问共用同一套样本口径与特征命名，避免在不同章节出现互相矛盾的数字。
